% Thesis skeleton in Chinese (Traditional characters in the text) with the ctex bundle. Compile with XeLaTeX:
%   xelatex thesis-skeleton-zh.tex   (run twice for the table of contents)
% ctexbook sets Chinese headings, spacing and font sizes (scheme = chinese, zihao = -4 below); the fonts are chosen by
% ctex's fontset detection for the operating system (fontset = mac, windows, ubuntu, fandol, ...). A university's own
% thesis class replaces this file when the school provides one; keep its rules over these.
% Every number in a thesis comes from an artifact (see ../../references/academic-papers-guide.md and
% scripts/render_results.py); [VALUE NEEDED] marks a value that has not been filled from one.
\documentclass[a4paper, zihao = -4, openany]{ctexbook}

\usepackage{amsmath, amssymb}
\usepackage{graphicx}
\usepackage{booktabs}
\usepackage{siunitx}
\usepackage[hidelinks]{hyperref}

\title{論文題目\\[0.5em]\large English Title of the Thesis}
\author{作者姓名}
\date{[VALUE NEEDED: 口試日期]}

\begin{document}

\frontmatter
\maketitle

\chapter*{摘要}
\addcontentsline{toc}{chapter}{摘要}
研究問題、方法、主要結果（含數值與不確定度：[VALUE NEEDED]）、結論。數值須與正文、表格一致。

\chapter*{Abstract}
\addcontentsline{toc}{chapter}{Abstract}
The English abstract states the same result with the same numbers ([VALUE NEEDED]) and the same claim strength.

\tableofcontents

\mainmatter

\chapter{緒論}
研究動機、物理背景與本論文的貢獻。

\chapter{實驗裝置與資料}
\section{探測器}
探測器與觸發；探測器細節引用其公開論文或實驗設定檔，不憑記憶。
\section{資料與模擬樣本}
資料期間、模擬樣本與版本。

\chapter{分析方法}
\section{事例選擇}
\section{背景估計}
\section{修正與系統不確定度}
系統不確定度表（表~\ref{tab:syst}）。

\begin{table}[htbp]
  \centering
  \caption{系統不確定度（[VALUE NEEDED]）}
  \label{tab:syst}
  \begin{tabular}{lc}
    \toprule
    來源 & 相對不確定度 \\
    \midrule
    {[VALUE NEEDED]} & {[VALUE NEEDED]} \\
    \bottomrule
  \end{tabular}
\end{table}

\chapter{結果}
測量結果為 $X = \text{[VALUE NEEDED]}$，單位以 \si{\giga\electronvolt} 等標準寫法表示。

\chapter{討論與結論}
結果的意義、限制與後續工作。

\backmatter
\begin{thebibliography}{9}
\bibitem{ref1} [作者], [題目], [期刊或 arXiv 編號] ([年份]).
\end{thebibliography}

\appendix
\chapter{補充材料}

\end{document}
